\section{弱收敛与弱星收敛}
\label{sec:lfunc-weak-convergence}

范数收敛要求整体误差$\lVert x_n-x\rVert$趋于零。在无限维空间中这个要求很强：第~\ref{chap:functions-as-vectors}章已经看到，单位球中的序列可以没有任何范数收敛的子列。换一个角度，只要求每一个固定的连续线性读数都收敛，会得到更弱、也更容易满足的收敛。本节分两种情形：在原空间$X$中用全部$\ell\in X^*$去读，得到弱收敛；在对偶空间$X^*$中只用原空间的向量去读，得到弱星收敛。两者的观测来源不同，不能只当作同一个概念的强弱两档。最后说明这两种收敛仍然控制序列的范数界。

\subsection{弱收敛}
\label{sec:lfunc-weak}

先用第~\ref{chap:topological-spaces}章的邻域语言描述“只看有限个读数时的接近”。给定$x_0\in X$、有限个泛函$\ell_1,\ldots,\ell_m\in X^*$和$\varepsilon>0$，令
\begin{equation}
U(x_0;\ell_1,\ldots,\ell_m;\varepsilon)=\{x\in X:\lvert\ell_i(x)-\ell_i(x_0)\rvert<\varepsilon,\ i=1,\ldots,m\}.
\label{eq:lfunc-weak-neighborhood}
\end{equation}
它由有限个读数的误差约束定义：只要这$m$个读数与$x_0$处的读数相差不到$\varepsilon$，就算“接近”。

\begin{definition}[弱拓扑与弱收敛]
\label{def:lfunc-weak-topology}
设$X$是赋范空间。若子集$G\subset X$对每个$x_0\in G$都包含某个形如式~\eqref{eq:lfunc-weak-neighborhood}的集合，称$G$为弱开集。弱开集组成的拓扑称为$X$上的\term{弱拓扑}（weak topology），记为$\sigma(X,X^*)$。序列$(x_n)$在弱拓扑中收敛到$x$，称为\term{弱收敛}（weak convergence），记为$x_n\rightharpoonup x$。
\end{definition}

需要核对弱开集确实组成拓扑。空集和$X$显然是弱开集，任意并的条件逐点检查即可。设$x_0$属于两个弱开集$G_1,G_2$之交，$G_1$含$U(x_0;\ell_1,\ldots,\ell_m;\varepsilon)$，$G_2$含$U(x_0;\ell'_1,\ldots,\ell'_{m'};\varepsilon')$；把两组泛函合在一起、误差取$\min(\varepsilon,\varepsilon')$，得到的集合同时含于两者，所以交仍是弱开集。式~\eqref{eq:lfunc-weak-neighborhood}本身也是弱开集：若$x_1$属于它，取$\delta=\varepsilon-\max_i\lvert\ell_i(x_1)-\ell_i(x_0)\rvert>0$，则$U(x_1;\ell_1,\ldots,\ell_m;\delta)$含于其中。

\begin{theorem}[弱收敛的读数判别]
\label{thm:lfunc-weak-criterion}
在赋范空间$X$中，$x_n\rightharpoonup x$当且仅当$\ell(x_n)\to\ell(x)$对每个$\ell\in X^*$成立。弱极限若存在则唯一。
\end{theorem}

\begin{proof}
设$x_n\rightharpoonup x$。对固定的$\ell$和$\varepsilon>0$，$U(x;\ell;\varepsilon)$是$x$的弱开邻域，序列最终落在其中，即$n$充分大时$\lvert\ell(x_n)-\ell(x)\rvert<\varepsilon$。反过来，设每个读数都收敛。任给$x$的弱开邻域，它包含某个$U(x;\ell_1,\ldots,\ell_m;\varepsilon)$；对每个$i$取$N_i$使$n\geq N_i$时$\lvert\ell_i(x_n)-\ell_i(x)\rvert<\varepsilon$，则$n\geq\max_iN_i$时$x_n$落在这个邻域中。这里只需对有限个$i$取最大值，所以不需要不同读数的收敛速度一致。

唯一性：若$x_n\rightharpoonup x$且$x_n\rightharpoonup y$，则$\ell(x)=\ell(y)$对一切$\ell\in X^*$成立，由定理~\ref{thm:lfunc-norming}得$x=y$。
\end{proof}

唯一性用到了连续线性读数足以区分向量，这是Hahn--Banach定理的结论；同样的理由说明弱拓扑是Hausdorff的。在Hilbert空间中，由Riesz表示定理，判别式变成$\langle y,x_n\rangle\to\langle y,x\rangle$对每个$y\in H$成立：弱收敛就是在每个固定方向上的内积收敛。

\begin{theorem}[范数收敛推出弱收敛]
\label{thm:lfunc-strong-implies-weak}
若$\lVert x_n-x\rVert\to0$，则$x_n\rightharpoonup x$。
\end{theorem}

\begin{proof}
对每个$\ell\in X^*$，$\lvert\ell(x_n)-\ell(x)\rvert\leq\lVert\ell\rVert\lVert x_n-x\rVert\to0$。
\end{proof}

\paragraph{逆向不成立}
在$\ell^2$中考察标准单位向量$\bm e_n$。两两距离为$\sqrt2$，没有范数收敛的子列。但对固定的$y\in\ell^2$，$\langle y,\bm e_n\rangle=\overline{y_n}$，由$\sum_n\lvert y_n\rvert^2<\infty$知它趋于零，所以$\bm e_n\rightharpoonup0$，而$\lVert\bm e_n\rVert=1$始终不变。更一般地，Hilbert空间中任何标准正交序列$(u_n)$都弱收敛到零，因为由Bessel不等式$\sum_n\lvert\langle y,u_n\rangle\rvert^2\leq\lVert y\rVert^2$，通项趋于零。

函数空间中的典型例子是振荡。$\{\sin(nx)/\sqrt\pi\}_{n\geq1}$是$L^2(0,2\pi)$中的标准正交序列，所以$\sin(nx)\rightharpoonup0$。取读数$g=\mathbf{1}_{[0,1]}$可以直接看到：$\int_0^1\sin(nx)\,dx=(1-\cos n)/n\to0$，而$\lVert\sin(nx)\rVert_2=\sqrt\pi$。函数在正负之间越摆越快，任何固定的权重与它相乘再积分，正负部分都相互抵消。弱收敛失去的正是对整体能量的控制：每个固定方向上的分量都在消失，能量却转移到越来越高的频率或越来越远的坐标上。

\paragraph{无限维中弱邻域很大}
弱邻域$U(0;\ell_1,\ldots,\ell_m;\varepsilon)$包含子空间$\ker\ell_1\cap\cdots\cap\ker\ell_m$，后者是$m$个线性方程的解集。若$X$是无限维的，映射$x\mapsto(\ell_1(x),\ldots,\ell_m(x))\in\mathbb F^m$不可能单射，这个交集含有非零向量，从而含有整条直线。所以无限维空间中每个弱邻域都是无界的，弱拓扑严格弱于范数拓扑。有限维时，坐标泛函有限个，弱邻域可以取成坐标误差都小于$\varepsilon$的集合，弱收敛与范数收敛一致。

\paragraph{部分读数不能代替全部读数}
弱收敛要求全部$\ell\in X^*$都收敛。只检查一部分读数，结论可能完全不同。在$\ell^2$中取$x_n=n\bm e_n$，每个坐标泛函$x\mapsto x_k$在$x_n$上的值最终为零，坐标逐个收敛到零向量。但取$y=(1/k)_k\in\ell^2$，$\langle y,x_n\rangle=1$不趋于零，所以$x_n$不弱收敛到零；由下面的定理~\ref{thm:lfunc-weak-bounded}，它无界，也不弱收敛到任何向量。坐标泛函的线性组合在$(\ell^2)^*$中稠密，可它们只是全部读数的一部分；缺少统一的范数界时，稠密子集上的收敛推不到全部读数上。

\subsection{弱星收敛}
\label{sec:lfunc-weak-star}

当变化的对象本身是读数，即$X^*$中的元素时，还有另一种自然的测试方法：只用原空间$X$中的向量去读它们。对$\ell_0\in X^*$、有限个向量$x_1,\ldots,x_m\in X$和$\varepsilon>0$，令
\[
U^*(\ell_0;x_1,\ldots,x_m;\varepsilon)=\{\ell\in X^*:\lvert\ell(x_i)-\ell_0(x_i)\rvert<\varepsilon,\ i=1,\ldots,m\}.
\]

\begin{definition}[弱星拓扑与弱星收敛]
\label{def:lfunc-weak-star}
设$X$是赋范空间。用上述集合代替式~\eqref{eq:lfunc-weak-neighborhood}，按定义~\ref{def:lfunc-weak-topology}的方式得到$X^*$上的拓扑，称为\term{弱星拓扑}（weak-star topology），记为$\sigma(X^*,X)$。序列$(\ell_n)$在其中收敛到$\ell$，称为\term{弱星收敛}（weak-star convergence），记为$\ell_n\overset{*}{\rightharpoonup}\ell$。
\end{definition}

与定理~\ref{thm:lfunc-weak-criterion}的证明完全相同，$\ell_n\overset{*}{\rightharpoonup}\ell$当且仅当$\ell_n(x)\to\ell(x)$对每个$x\in X$成立，即逐个向量的求值收敛。弱星极限唯一：若$\ell(x)=\ell'(x)$对一切$x$成立，则$\ell=\ell'$，这里不需要延拓定理。

\paragraph{与对偶空间自身的弱收敛比较}
$X^*$本身是赋范空间，它也有定义~\ref{def:lfunc-weak-topology}意义下的弱拓扑$\sigma(X^*,X^{**})$，测试对象是$X^{**}$的全部元素。每个$x\in X$通过$\Phi_x(\ell)=\ell(x)$给出$X^{**}$的一个元素（第~\ref{sec:lfunc-reflexivity}小节将它记为$J(x)$），所以弱星收敛只用$X^{**}$中这一部分元素测试，要求更少：
\[
\ell_n\rightharpoonup\ell\ \text{（在}X^*\text{中）}\quad\Longrightarrow\quad\ell_n\overset{*}{\rightharpoonup}\ell .
\]
逆向一般不成立，因为$X^{**}$中可以有不来自$X$的元素。

\begin{example}[弱星收敛而不弱收敛]
\label{ex:lfunc-weak-star-not-weak}
取$X=\ell^1$。由定理~\ref{thm:lfunc-lp-dual}（计数测度，$p=1$），$X^*$可以等同于$\ell^\infty$，$y\in\ell^\infty$读作$x\mapsto\sum_ky_kx_k$。令$y^{(n)}\in\ell^\infty$在第$n$项及以后为$1$、此前为$0$。对每个$x\in\ell^1$，$\sum_ky^{(n)}_kx_k=\sum_{k\geq n}x_k\to0$，所以$y^{(n)}\overset{*}{\rightharpoonup}0$。

但$y^{(n)}$在$\ell^\infty$中不弱收敛。取第~\ref{sec:lfunc-hahn-banach}小节构造的$\Lambda\in(\ell^\infty)^*$，它在收敛数列上取极限，于是$\Lambda(y^{(n)})=1$对一切$n$成立，不趋于$\Lambda(0)=0$。若$y^{(n)}$弱收敛到某个$y$，由上面的蕴含它也弱星收敛到$y$，由弱星极限唯一得$y=0$，矛盾。$\Lambda$正是$X^{**}$中不来自$X$的元素。
\end{example}

\subsection{弱收敛列的范数控制}
\label{sec:lfunc-weak-norm-control}

弱收敛只给出逐个读数的信息，但借助一致有界原理，它仍然控制了序列的范数。关键的想法是把向量看成对偶空间上的读数：$x\in X$对应$X^*$上的泛函$\ell\mapsto\ell(x)$，由式~\eqref{eq:lfunc-norm-dual}，这个泛函的范数恰为$\lVert x\rVert$。

\begin{theorem}[弱收敛列有界且范数下半连续]
\label{thm:lfunc-weak-bounded}
设$X$是赋范空间，$x_n\rightharpoonup x$。则$\sup_n\lVert x_n\rVert<\infty$，并且
\begin{equation}
\lVert x\rVert\leq\liminf_{n\to\infty}\lVert x_n\rVert .
\label{eq:lfunc-weak-lsc}
\end{equation}
\end{theorem}

\begin{proof}
\emph{有界性。}令$\Phi_n(\ell)=\ell(x_n)$，$\ell\in X^*$。$\Phi_n$是$X^*$上的有界线性泛函，由式~\eqref{eq:lfunc-norm-dual}，$\lVert\Phi_n\rVert=\lVert x_n\rVert$。对每个固定的$\ell$，$\Phi_n(\ell)=\ell(x_n)$收敛，因而有界，所以$\{\Phi_n\}$在$X^*$上逐点有界。由定理~\ref{thm:lfunc-dual-complete}，$X^*$是Banach空间，由定理~\ref{thm:lfunc-uniform-boundedness}得$\sup_n\lVert x_n\rVert=\sup_n\lVert\Phi_n\rVert<\infty$。

\emph{下半连续。}任取$\lVert\ell\rVert\leq1$，则$\lvert\ell(x)\rvert=\lim_n\lvert\ell(x_n)\rvert\leq\liminf_n\lVert x_n\rVert$。对$\ell$取最大值，由式~\eqref{eq:lfunc-norm-dual}得式~\eqref{eq:lfunc-weak-lsc}。
\end{proof}

证明中用到完备性的空间是$X^*$而不是$X$，所以结论对不完备的$X$也成立。式~\eqref{eq:lfunc-weak-lsc}可以是严格不等式：$\bm e_n\rightharpoonup0$，而$\lVert0\rVert=0<1$。弱极限允许一部分能量在极限中消失，但不会凭空增加。

\paragraph{弱星收敛列}
对弱星收敛，逐点有界直接发生在$X$上。若$X$是Banach空间且$\ell_n\overset{*}{\rightharpoonup}\ell$，则对每个$x$，$(\ell_n(x))$收敛因而有界，由定理~\ref{thm:lfunc-uniform-boundedness}得$\sup_n\lVert\ell_n\rVert<\infty$；同样有$\lVert\ell\rVert\leq\liminf_n\lVert\ell_n\rVert$，因为$\lvert\ell(x)\rvert=\lim\lvert\ell_n(x)\rvert\leq\liminf_n\lVert\ell_n\rVert$对单位向量$x$成立。这里$X$的完备性不能省：在第~\ref{sec:lfunc-uniform-boundedness}节的$c_{00}$中，$\ell_n(x)=nx_n$对每个$x$最终为零，所以$\ell_n\overset{*}{\rightharpoonup}0$，但$\lVert\ell_n\rVert=n$无界。

\paragraph{有界只是必要条件}
定理~\ref{thm:lfunc-weak-bounded}给出了一个检验：无界序列不可能弱收敛。例如$\sqrt n\,\bm e_n$在$\ell^2$中不弱收敛，可以直接取$y_k=k^{-3/4}$验证，$y\in\ell^2$而$\langle y,\sqrt n\,\bm e_n\rangle=n^{1/4}\to\infty$。但有界不能反推收敛：$x_n=(-1)^n\bm e_1$有界，在读数$\langle\bm e_1,\cdot\rangle$下交替取$\pm1$，整列不弱收敛。有界序列是否至少有弱收敛的子列，是下一节的紧性问题，答案取决于空间。
